%==============================================================
% CV Jorge Armando Cordero García
% VERSIÓN 2: Ingeniero de Infraestructura IT / Redes / Sistemas
% Optimizado para ATS | MiKTeX | VS Code
% Compilar con: pdflatex CV_Jorge_Cordero_Infraestructura.tex
%==============================================================

\documentclass[10pt, letterpaper]{article}

%------ PAQUETES ------
\usepackage[utf8]{inputenc}
\usepackage[T1]{fontenc}
\usepackage{babel}
\usepackage[left=1.8cm, right=1.8cm, top=1.5cm, bottom=1.5cm]{geometry}
\usepackage{lmodern}
% \usepackage{microtype}
\usepackage{parskip}
\usepackage{enumitem}
% \usepackage{titlesec}
\usepackage{xcolor}
\usepackage{hyperref}
\usepackage{amssymb}
% \usepackage{fontawesome5}
\usepackage{multicol}
\usepackage{array}
\usepackage{tabularx}
% \\usepackage{colortbl}
\usepackage{tikz}

%------ COLORES ------
\definecolor{primary}{HTML}{0D3349}
\definecolor{accent}{HTML}{1B6CA8}
\definecolor{accentlight}{HTML}{E8F4FD}
\definecolor{lightgray}{HTML}{F7F7F7}
\definecolor{darkgray}{HTML}{333333}
\definecolor{medgray}{HTML}{555555}
\definecolor{rulecol}{HTML}{1B6CA8}
\definecolor{tagbg}{HTML}{E3EFF8}

%------ HYPERLINKS ------
\hypersetup{
    colorlinks=true,
    urlcolor=accent,
    linkcolor=primary,
    pdfauthor={Jorge Armando Cordero García},
    pdftitle={CV - Ingeniero de Infraestructura IT},
    pdfsubject={Curriculum Vitae}
}

%------ FORMATO DE SECCIONES ------
% Formato de sección simplificado para el entorno actual
% Se usa el estilo por defecto de sección

%------ COMANDOS PERSONALIZADOS ------
\newcommand{\cvheader}[6]{%
  \noindent
  \begin{minipage}[t]{0.62\linewidth}
    {\LARGE\bfseries\color{primary} #1}\\[4pt]
    {\large\color{accent}\bfseries #2}
  \end{minipage}%
  \hfill
  \begin{minipage}[t]{0.36\linewidth}
    \raggedleft
    {\small\color{darkgray}
      Ubicaci\u00f3n: #3\\[2pt]
      Tel: #4\\[2pt]
      Email: \href{mailto:#5}{#5}\\[2pt]
      LinkedIn: \href{https://linkedin.com/in/jacoga}{linkedin.com/in/jacoga}
    }
  \end{minipage}
  \vspace{4pt}
  \textcolor{rulecol}{\rule{\linewidth}{2pt}}
  \vspace{2pt}
}

\newcommand{\cventry}[5]{%
  \noindent
  \begin{tabularx}{\linewidth}{@{}X r@{}}
    {\bfseries\color{primary} #1} & \cellcolor{accentlight}{\small\color{accent}\bfseries #2} \\
    {\small\color{accent} #3} \textcolor{medgray}{$|$} {\small\color{darkgray}\itshape #4} & \\
  \end{tabularx}
  \vspace{2pt}
  #5
  \vspace{6pt}
}

\newcommand{\keyword}[1]{\textbf{#1}}

\setlist[itemize]{
  leftmargin=14pt,
  itemsep=1.5pt,
  parsep=0pt,
  topsep=2pt,
  label=\textcolor{accent}{$\circ$}
}

\newcommand{\tag}[1]{%
  \tikz[baseline=(t.base)]\node[fill=tagbg, rounded corners=2pt, inner sep=3pt, text=accent](t){\footnotesize\bfseries #1};%
}

\pagestyle{empty}

%==============================================================
\begin{document}
%==============================================================

% ENCABEZADO
\cvheader
  {Jorge Armando Cordero García}
  {Ingeniero de Sistemas | Infraestructura IT \& Redes}
  {Gamarra, Cesar – Colombia}
  {+57 317 819 5354}
  {iscuande\_1982@hotmail.com}
  {linkedin.com/in/jacoga}

%------ PERFIL PROFESIONAL ------
\section{Perfil Profesional}

Ingeniero de Sistemas con más de \textbf{10 años de experiencia} en administración y gestión de infraestructura tecnológica \emph{on-premise} en entidades del sector público, privado, salud y educación. Especializado en \textbf{administración de servidores Windows Server}, \textbf{redes LAN/WAN}, \textbf{Active Directory}, cableado estructurado y soporte técnico multinivel. Experiencia probada garantizando hasta \textbf{99,5\% de disponibilidad operativa} en entornos críticos. Formación complementaria en ciberseguridad, desarrollo web frontend e inteligencia artificial aplicada. Certificado por \textbf{Google} y \textbf{Cisco Networking Academy}. Disponibilidad inmediata, modalidad presencial/remota/híbrida, con movilidad nacional.

%------ EXPERIENCIA PROFESIONAL ------
\section{Experiencia Profesional}

\cventry
  {Ingeniero de Sistemas – Consultor TI}
  {Ene 2024 – Presente}
  {Construcciones \& Soluciones del Río S.A.S. / Tictelco S.A.S. / Clientes Privados}
  {Aguachica, Cesar – Remota/Presencial}
  {
    \begin{itemize}
      \item Consultoría integral de infraestructura IT para empresas del sector construcción, salud y PYME.
      \item Administración de servidores locales, gestión de racks, switching y diagnóstico de conectividad WAN/LAN.
      \item Mantenimiento preventivo y correctivo de equipos de cómputo e impresoras; instalación y configuración de equipos.
      \item Revisión y certificación de red cableada estructurada; mantenimiento de cableado físico.
      \item Implementación de soluciones tecnológicas para mejora de procesos operativos.
      \item Desarrollo y mantenimiento web frontend: \keyword{HTML5, CSS3, JavaScript, WordPress}.
      \item Capacitación técnica a usuarios finales y transferencia de conocimiento tecnológico.
    \end{itemize}
  }

\cventry
  {Ingeniero de Sistemas}
  {Feb 2023 – Jul 2024}
  {E.S.E. Hospital Olaya Herrera}
  {Gamarra, Cesar – Presencial}
  {
    \begin{itemize}
      \item Administración integral del sistema de información asistencial hospitalaria (SISPRO, MIPRES, RUAF, SIGIRES).
      \item Gestión de \keyword{servidores Windows Server} y red LAN; \textbf{99,5\% de disponibilidad} en entorno crítico.
      \item Soporte técnico especializado a más de \textbf{200 usuarios} con resolución eficiente de incidencias.
      \item Mantenimiento de infraestructura de red LAN, equipos de cómputo, impresoras y activos TI.
      \item Gestión del inventario y ciclo de vida de activos tecnológicos institucionales.
    \end{itemize}
  }

\cventry
  {Ingeniero de Sistemas}
  {Feb 2022 – Jul 2022}
  {Universidad Popular del César}
  {Aguachica, Cesar – Presencial}
  {
    \begin{itemize}
      \item Administración de infraestructura tecnológica educativa: servidores, redes LAN y equipos de cómputo.
      \item Control y gestión de inventario TI; administración del ciclo de vida de activos tecnológicos.
      \item Liderazgo en proyectos de incorporación de TIC en procesos académicos institucionales.
      \item Capacitación docente y administrativa en herramientas digitales y plataformas educativas.
    \end{itemize}
  }

\cventry
  {Líder GELT – Gestión TIC Municipal}
  {Jun 2015 – Feb 2016}
  {Alcaldía Municipal de Gamarra}
  {Gamarra, Cesar – Presencial}
  {
    \begin{itemize}
      \item Liderazgo del área GELT (Gestión Electrónica) municipal; planeación y ejecución de proyectos TIC.
      \item Administración de servidores locales y red LAN institucional; gestión de recursos e inventario TIC.
      \item Articulación interinstitucional para implementación de soluciones de gobierno electrónico.
    \end{itemize}
  }

\cventry
  {Ingeniero de Sistemas}
  {Mar 2016 – Mar 2017}
  {Linktech de Colombia}
  {Bogotá, D.C. – Presencial}
  {
    \begin{itemize}
      \item Mantenimiento preventivo/correctivo de equipos; instalación y asesoría en hardware/software.
      \item Administración de redes, cableado estructurado y certificación de infraestructura de red corporativa.
    \end{itemize}
  }

\cventry
  {Ingeniero de Sistemas}
  {Ene 2017 – Dic 2017}
  {Cordisco}
  {Bogotá, D.C. – Presencial}
  {
    \begin{itemize}
      \item Administración de red corporativa y soporte técnico interno; implementación de mejoras tecnológicas.
    \end{itemize}
  }

%------ FORMACIÓN ACADÉMICA ------
\section{Formación Académica}

\noindent
\begin{tabularx}{\linewidth}{@{}X r@{}}
  {\bfseries\color{primary} Ingeniería de Sistemas} & {\small\color{medgray} 2015 – 2023} \\
  {\small\color{darkgray} Universidad Popular del César – Aguachica, Cesar} & \\[6pt]
  {\bfseries\color{primary} Especialización en TIC para Diseño de Estrategias Didácticas} & {\small\color{medgray} 2022 – 2023} \\
  {\small\color{darkgray} Universidad Popular del César – En proceso de grado} & \\
\end{tabularx}

%------ CERTIFICACIONES ------
\section{Certificaciones Relevantes}

\begin{multicols}{2}
\begin{itemize}
  \item \textbf{Google IT Support Professional} (2025)
  \item \textbf{Google Data Analytics} (2025)
  \item \textbf{Google Cybersecurity} – En curso (2025)
  \item \textbf{Cisco: IT Customer Support Basics} (2025)
  \item \textbf{Cisco: Introduction to IoT} (2025)
  \item \textbf{Cisco: Introduction to Cybersecurity} (2025)
  \item \textbf{Fundamentos de IA – Google} (2024)
  \item \textbf{IA Generativa – Microsoft} (2024)
  \item \textbf{HTML5/CSS3, JS ES6+, React.js} – LinkedIn Learning (2024)
  \item \textbf{Inglés Técnico Niveles 1–3} – SENA (2021)
\end{itemize}
\end{multicols}

%------ HABILIDADES TÉCNICAS ------
\section{Habilidades Técnicas}

\vspace{4pt}
\noindent
\tag{Windows Server} \tag{Active Directory} \tag{Linux} \tag{TCP/IP} \tag{LAN/WAN} \tag{Switching} \tag{Routing} \tag{Cableado Estructurado} \tag{Rack Management} \tag{ITIL} \tag{Help Desk} \tag{TeamViewer} \tag{AnyDesk} \tag{HTML5} \tag{CSS3} \tag{JavaScript} \tag{WordPress} \tag{React.js} \tag{Git} \tag{MS Office} \tag{Google Workspace} \tag{Ciberseguridad} \tag{Virtualización}

\vspace{6pt}

\noindent\textbf{\color{primary}Gestión:} Administración de proyectos IT $\cdot$ Control de inventario tecnológico $\cdot$ Documentación técnica $\cdot$ Gestión de activos TI $\cdot$ Continuidad operativa

%------ IDIOMAS ------
\section{Idiomas}

\noindent
\begin{tabularx}{\linewidth}{@{}l X@{}}
  \textbf{Español:} & Nativo \\
  \textbf{Inglés:} & Básico–Técnico (A2) – Lectura de documentación técnica en inglés \\
\end{tabularx}

%------ PIE ------
\vspace{8pt}
\begin{center}
\textcolor{medgray}{\small \faCheckCircle[regular]~Disponibilidad inmediata \quad \faCheckCircle[regular]~Presencial / Remota / Híbrida \quad \faCheckCircle[regular]~Movilidad nacional}
\end{center}

\end{document}
